\section{Universal constructions with explicit operations}\label{sec:limits}

In this section, an equality of parallel arrows always means the chosen
hom-setoid equality.  We suppress its source and target parameters, but never
identify two object codes merely because their objects are isomorphic.  Every
operation mentioned below is a fixed primitive recursive operation on codes.
Its typing and algebraic laws are hypotheses already proved for the
presentation under consideration.  An operation may return arbitrary junk on
inputs outside its stated domain.

\subsection{Terminal objects, products, and uniqueness}

A \emph{chosen terminal object} consists of an object $1$ and an operation
$!_a:a\to1$ such that every arrow $u:a\to1$ equals $!_a$.  A chosen initial
object is the dual construction.  For an initial object $0$, its unique arrow
$0\to a$ is supplied as an operation; the statement that an arrow exists for
each $a$ is not silently converted into such an operation.

A \emph{chosen binary product} of $a,b$ consists of an object $a\times b$,
projections $p:a\times b\to a$ and $q:a\times b\to b$, and a pairing operation
\[
 (u:x\to a,\ v:x\to b)\longmapsto\langle u,v\rangle:x\to a\times b
\]
satisfying
\[
 p\langle u,v\rangle=u,\qquad q\langle u,v\rangle=v,\qquad
 \langle pw,qw\rangle=w.
\]
Composition is written by juxtaposition, with the arrow on the right applied
first.  These three equations include both existence and uniqueness.  They
also prove that pairing respects equality: if $u=u'$ and $v=v'$, the two
pairings have equal projections and hence are equal.

\begin{proposition}[Uniqueness of universal objects]\label{prop:universal-unique}
Two terminal objects are uniquely isomorphic.  Two products of the same
ordered pair are uniquely isomorphic by an isomorphism preserving projections.
More generally, two limits of a fixed diagram are uniquely isomorphic by an
isomorphism preserving their cones.  All these isomorphisms and their inverses
are obtained from the supplied universal operations by composition.
\end{proposition}
\begin{proof}
For terminal objects $t,t'$, let $u:t\to t'$ and $v:t'\to t$ be their supplied
unique arrows.  Since $vu$ and $\id_t$ both map $t$ to $t$, uniqueness at $t$
gives $vu=\id_t$; similarly $uv=\id_{t'}$.  Any arrow $t\to t'$ equals $u$.

For products $(p:P\to a,q:P\to b)$ and $(p':P'\to a,q':P'\to b)$, put
$u=\langle p,q\rangle':P\to P'$ and
$v=\langle p',q'\rangle:P'\to P$.  Then $pvu=p'u=p$ and $qvu=q'u=q$, so
product uniqueness gives $vu=\id_P$.  The other composite is analogous.
An arrow preserving both projections must equal the indicated pairing.

For limits, let $u:L\to L'$ be the mediating arrow from the cone of $L$
into the limiting cone of $L'$, and define $v:L'\to L$ conversely.  Every
projection of $vu$ equals the corresponding projection of $\id_L$; limit
uniqueness gives their equality.  Interchanging the two cones proves the
other composite equation and the same uniqueness argument proves the final
claim.  These are finite derivations in intuitionistic first-order logic:
neither excluded middle nor induction is used.
\end{proof}

An \emph{equalizer} of $f,g:a\to b$ is an object $E$ and $e:E\to a$ with
$fe=ge$, together with a supplied operation taking $u:x\to a$ and the premise
$fu=gu$ to an arrow $\operatorname{lift}(u):x\to E$.  Its laws are
\[
 e\operatorname{lift}(u)=u,
 \qquad ev=u\ \Longrightarrow\ v=\operatorname{lift}(u).
\]
The operation need not test the premise: its correctness is asserted only
under that premise.  Equalizers are monomorphisms: if $er=es$, both $r$ and
$s$ are mediating arrows for the same arrow $er$, so they are equal.

A \emph{pullback} of $f:a\to c$ and $g:b\to c$ consists of
$p:P\to a$, $q:P\to b$ with $fp=gq$, and a supplied operation
$(u,v)\mapsto\langle u,v\rangle_{f,g}$ for compatible pairs $fu=gv$, having
those projections and uniquely determined by them.

\begin{proposition}[Pullbacks from products and equalizers]
If the relevant chosen product and equalizer are supplied, the pullback of
$f:a\to c$ and $g:b\to c$ is obtained by equalizing
\[
 f p_a,\ g p_b:a\times b\longrightarrow c.
\]
Conversely, a terminal object and chosen pullbacks produce binary products;
chosen binary products and pullbacks produce equalizers.
\end{proposition}
\begin{proof}
Let $e:P\to a\times b$ equalize $fp_a,gp_b$ and give $P$ the projections
$p_a e,p_b e$.  Their compatibility is the equalizer equation.  If $fu=gv$,
then $\langle u,v\rangle$ equalizes the two arrows, so its equalizer lift has
the required projections.  An arrow with those projections satisfies
$ew=\langle u,v\rangle$ by product uniqueness; equalizer uniqueness then
identifies it with the lift.

The pullback of the two unique arrows $a\to1\leftarrow b$ is a product,
because every pair $u:x\to a,v:x\to b$ is compatible by terminal uniqueness.
For the last assertion, pull back $(f,g):a\to b\times b$ along the diagonal
$\Delta_b=\langle\id_b,\id_b\rangle:b\to b\times b$.  Write the resulting
projections as $e:E\to a$ and $r:E\to b$.  The pullback equation gives
$fe=r=ge$.  If $fu=gu$, the pair $(u,fu)$ is compatible and therefore has a
unique lift $w$ with $ew=u$.  Any other $w'$ with $ew'=u$ has
$rw'=few'=fu$, so pullback uniqueness proves $w'=w$.
\end{proof}

\subsection{All fixed finite limits}

Fix externally a finite category $J$: a finite list of objects, a finite list
of arrows with specified endpoints, and fixed identity and composition
tables.  For a diagram $D:J\to C$, a cone with vertex $x$ is a finite tuple
$c_i:x\to D_i$ satisfying $D(u)c_i=c_j$ for every listed arrow $u:i\to j$.
A chosen limit consists of a cone $p_i:L\to D_i$ and a primitive recursive
operation $m(x,(c_i))$ such that $p_i m=c_i$, and such that any arrow with
those equations equals $m$.  A cone predicate is an arithmetical formula; it
need not be decidable.  The operation is only required to have its stated
properties under that predicate.

\begin{theorem}[Finite-limit construction]\label{thm:finite-limits}
For each fixed finite category $J$, chosen terminal objects, binary products,
and equalizers yield a chosen limit for every $J$-diagram.  The constructor
and the mediating operation are primitive recursive compositions of the
supplied operations.  The proof translates into \ipra{} without any induction
on a formula containing the category's equality predicate.
\end{theorem}
\begin{proof}
Form the finite products
\[
 A=\prod_{i\in\operatorname{Ob}J}D_i,
 \qquad B=\prod_{(u:i\to j)\in\operatorname{Arr}J}D_j.
\]
They mean fixed parenthesized iterations of binary products; an empty product
means the terminal object.  Write $\pi_i:A\to D_i$ and
$\rho_u:B\to D_j$.  By repeated pairing define $s,t:A\to B$ by
\[
 \rho_u s=D(u)\pi_i,\qquad \rho_u t=\pi_j.
\]
Take their equalizer $e:L\to A$ and put $p_i=\pi_i e$.  Then
$D(u)p_i=D(u)\pi_i e=\rho_u se=\rho_u te=\pi_j e=p_j$, so these arrows
form a cone.

Given a cone $(c_i)$, repeated pairing supplies $c:x\to A$ with
$\pi_i c=c_i$.  For every arrow $u:i\to j$,
$\rho_u sc=D(u)c_i=c_j=\rho_u tc$.  Product uniqueness gives $sc=tc$.
Consequently its equalizer lift $m:x\to L$ exists by the supplied operation
and satisfies $p_i m=\pi_i em=\pi_i c=c_i$.  If $v:x\to L$ has the same
cone projections, then all product projections of $ev$ and $c$ agree, so
$ev=c$.  Equalizer uniqueness yields $v=m$.

All lists in this proof have an externally fixed length.  Thus each use of
``for every listed arrow'' expands to a finite conjunction, and repeated
pairing expands to a fixed term.  The argument is a finite intuitionistic
proof for that shape.  It does not infer a single internal theorem for
arbitrarily long diagrams by induction on an undecidable predicate.  A
uniform varying-length version requires an additional verified coding and
invariant argument, as discussed in the foundational section.
\end{proof}

Applying the same argument in $C^{\op}$ proves the dual theorem: a chosen
initial object, binary coproducts, and coequalizers give all fixed finite
colimits.  In particular, a pushout is the coequalizer of the two induced
arrows into a coproduct.  The hypotheses supply the coequalizer operation;
this theorem does not manufacture quotient objects in an arbitrary category.

\subsection{Idempotents and their canonical completion}

An idempotent is an endomorphism $e:a\to a$ satisfying $e^2=e$.  A
\emph{splitting} consists of $r:a\to b$ and $s:b\to a$ satisfying
$sr=e$ and $rs=\id_b$.  A splitting is therefore a displayed retraction,
not an unspecified existential witness.

\begin{theorem}[Explicit idempotent completion]\label{thm:karoubi}
Every presented category $C$ has a presented category $\operatorname{Kar}(C)$
whose objects are pairs $(a,e)$ with $e^2=e$, and whose arrows
$(a,e)\to(b,d)$ are arrows $f:a\to b$ satisfying
\[
 fe=f=df.
\]
Composition is inherited from $C$, and the identity of $(a,e)$ is $e$.
The inclusion $a\mapsto(a,\id_a)$ is fully faithful with the identity
operation on hom-codes.  Every idempotent in $\operatorname{Kar}(C)$ has a
canonical primitive recursive splitting.
\end{theorem}
\begin{proof}
If $f:(a,e)\to(b,d)$ and $g:(b,d)\to(c,k)$, then
$(gf)e=g(fe)=gf$ and $k(gf)=(kg)f=gf$, so $gf$ is a valid arrow.
The equations $fe=f=df$ make $e,d$ the appropriate identities;
associativity follows from that in $C$.  For identity idempotents the support
conditions impose no extra restriction, proving full faithfulness.

Let $p:(a,e)\to(a,e)$ satisfy $p^2=p$.  The object $(a,p)$ is valid.
Use the underlying arrow $p$ for both
$r:(a,e)\to(a,p)$ and $s:(a,p)\to(a,e)$.  Their support equations follow
from $pe=p=ep$ and $p^2=p$.  We have $sr=p$ as an endomorphism of $(a,e)$,
while $rs=p=\id_{(a,p)}$.  Pairing and the projection operations on codes
provide the asserted constructors.
\end{proof}

For completeness, suppose $F:C\to D$ is supplied and that every $F(e)$ has
a chosen splitting $F(a)\xrightarrow{r_e}B_e\xrightarrow{s_e}F(a)$,
primitive recursively in the displayed idempotent.  Then
\[
 \widehat F(a,e)=B_e,\qquad
 \widehat F(f)=r_d F(f)s_e
\]
defines a functor.  Indeed,
$r_eF(e)s_e=r_es_er_es_e=\id_{B_e}$, and
\[
 (r_kF(g)s_d)(r_dF(f)s_e)
 =r_kF(g)F(d)F(f)s_e
 =r_kF(gf)s_e.
\]
For the inclusion $I:C\to\operatorname{Kar}(C)$, this is an extension
\emph{up to a specified natural isomorphism}, not necessarily a literal
identity of object codes.  Its component
$s_{\id_a}:\widehat F(Ia)\to Fa$ has inverse $r_{\id_a}$, and
\[
 s_{\id_b}\widehat F(I f)
 =s_{\id_b}r_{\id_b}F(f)s_{\id_a}
 =F(f)s_{\id_a}.
\]
Here $s_{\id_b}r_{\id_b}=F(\id_b)=\id_{Fb}$ modulo the given hom equality.
If the identity splittings are chosen as the canonical trivial splittings,
the object maps agree literally.  The theorem does not assert that arbitrary
idempotents already split in $C$ or $D$.
